% Source: https://artofproblemsolving.com/community/c3h1774223_prove_this_please
% Stable identity: aops-topic-1774223-q01
Let $M$ and $K$ be the midpoints of $AB$ and $BC$, and draw the perpendicular bisectors of $AB$ and $BC$ (the lines through $M$ and $K$ perpendicular to these sides).

\begin{asy}
usepackage("amsmath");
size(6cm);
defaultpen(fontsize(10pt));
pair A=(1,3.2),B=(0,0),C=(5,0);
pair cc(pair P,pair Q,pair R){pair a=Q-P,c=R-P; real d=2*(a.x*c.y-a.y*c.x);
  return P+((c.y*dot(a,a)-a.y*dot(c,c))/d,(a.x*dot(c,c)-c.x*dot(a,a))/d);}
pair O=cc(A,B,C); real R=abs(A-O);
pair M=(A+B)/2,K=(B+C)/2;
draw(circle(O,R),blue+linewidth(0.7));
draw(A--B--C--cycle,linewidth(0.9));
pair u1=unit(rotate(90)*(B-A)); draw(M-0.9*u1--O+0.9*unit(O-M),dashed);
draw(K+(0,-0.9)--O+(0,0.9),dashed);
draw(O--A,gray);draw(O--B,gray);draw(O--C,gray);
real s=0.22; pair e1=unit(A-B),e2=unit(O-M);
draw(M+s*e1--M+s*e1+s*e2--M+s*e2);
draw(K+(s,0)--K+(s,s)--K+(0,s));
dot(O);dot(M);dot(K);
label("$A$",A,N);label("$B$",B,SW);label("$C$",C,SE);label("$O$",O+0.32*dir(40));
label("$M$",M+0.4*dir(117.6));label("$K$",K,dir(-45));
label("\mbox{}",(min(currentpicture,user=true)+max(currentpicture,user=true))/2,black); // keeps the text after the figure black when LaTeX typesets the labels inline
\end{asy}

\textbf{The bisectors meet.} If they were parallel, then $AB$ and $BC$, being perpendicular to parallel lines, would be parallel; as they share the point $B$, the points $A$, $B$, $C$ would be collinear, which is not a triangle. So the two perpendicular bisectors meet at a point $O$.

\textbf{$O$ is equidistant from the vertices.} A point is on the perpendicular bisector of a segment exactly when it is equally far from the two endpoints (the right triangles $OMA$ and $OMB$ are congruent by SAS). So $OA=OB$ because $O$ is on the bisector of $AB$, and $OB=OC$ because $O$ is on the bisector of $BC$.

Hence $OA=OB=OC$, and the circle with centre $O$ and radius $OA$ passes through $A$, $B$ and $C$. Every triangle is cyclic. $\blacksquare$ (Moreover $OA=OC$ shows that $O$ also lies on the perpendicular bisector of $CA$, and the circle is unique, since its centre must lie on all three bisectors.)
